\appendix
\renewcommand{\thesubsection}{\thesection.\alph{subsection}\char41}
\numberwithin{equation}{section}
\renewcommand{\theequation}{\thesection.~\arabic{equation}}
\setlength{\parskip}{2mm}
\section{Representations of quantum mechanics}\label{representations}

\textbf{Complement of quantum mechanics}
An average is defined in quantum mechanics by:
\begin{align*}
    \expval{O}_\psi&=\expval{O}{\psi}\\
    \expval{O}_\rho&=\Tr{\rho O}
\end{align*}
We can see the two definitions are strictly equivalent, to illustrate let us pick a pure state $\rho=\ketbra{\psi}$ and the base of an Hamiltonian $\qty{\phi_n}_{n\in I_\mathcal{H}}$:
\begin{align*}
    \expval{O}_\rho&=\Tr{\rho O}\\
    &= \Tr{\ketbra{\psi} O}\\
    &= \Tr{\expval{O}{\psi}}
\end{align*}
However,\[\expval{O}{\psi}\in\mathbb{C}\]
Thus,
\begin{equation*}
    \expval{O}_\rho = \expval{O}{\psi}
\end{equation*}
\hfill $\square$

\subsection{Schrödinger and Heisenberg pictures}
In this appendix we will work by denoting the pictures with their acronyms, but we chose to lighten the notation of the total state $\rho_T\equiv\rho$ here. Please pay attention not to misinterpret $S$ in this part means \textit{in the Schrödinger picture}, whereas, $S$ in the main part means \textit{state of the system S}.\par

In the \textit{Schrödinger} picture we saw that $\ket{\psi(t)}=U(t)\ket{\psi}$ where, $U(t)=e^{-\frac{i}{\hbar}Ht}$ in the case of a time-independent Hamiltonian with $t_0=0$. The following demonstration can be made in the general case where the Hamiltonian is time dependent using Eq.\ (\ref{Htimedep}).\par

If we consider the average value of an observable $A_S(t)$ in the \textit{Schrödinger} picture,
\begin{align}
    \expval{A_S(t)} &= \expval{A_S}{\psi(t)}\nonumber\\
    &= \expval{\underbrace{U^\dagger(t)A_S(t)U(t)}_{A_H(t)}}{\psi}
\end{align}

This naturally let appear the \textit{Heisenberg representation} of observables and states.\\
Indeed, we have $\forall t,\;\ket{\psi(t)}_H=\ket{\psi(t_0)}_S\equiv\ket{\psi}_H$ which links both pictures at the initial time.\\
We can link both representation by calculating the evolution of observable in the \textit{Heisenberg picture}:
\begin{align*}
    \dert{A_H} &= \dert{}\left[U^\dagger A_{S}U\right]\\
    &= \dert{U^\dagger}A_{S}U + U^\dagger\derpart{A_S}U + \dert{U}A_{S}U\\
    &= \frac{i}{\hbar}U^\dagger H_{S}A_{S}U + U^\dagger\derpart{A}U - \frac{i}{\hbar}U^\dagger A_{S}H_{S}U\\
    &= \frac{i}{\hbar}[U^\dagger H_S\overbrace{UU^\dagger}^{\mathbb{I}} A_{S}U-U^\dagger A_{S}\overbrace{UU^\dagger}^{\mathbb{I}} H_{S}U]+ \derpart{A_S}\\
    &= \frac{i}{\hbar}\left[H_{H}A_{H}-A_{H}H_{H}\right] + \derpart{A_S}
\end{align*}
So,
\begin{equation}
    \boxed{\dert{A_H} = \frac{i}{\hbar}\left[H_{H},A_{H}\right] + \derpart{A_S}\label{eq:HSlink}}
\end{equation}

This link can even lead to the famous \textit{Ehrenfest theorem} for quantum mechanics:
\begin{equation}
    \boxed{\dert{\expval{A}}=\frac{i}{\hbar}\expval{\qty[H(t),A(t)]}+\expval{\derpart{A}}}
\end{equation}
We observe that if the observable doesn't depend explicitly on time (thus is constant), then the dynamics of the operator is fully describe by the \textit{Heisenberg representation}. 

\subsection{Interaction picture}
As mentioned in the main part, see~\ref{OQS}, we use perturbation theory to describe the interaction between the system $S$ and the environment $E$:
\begin{align}
    H_T &= \overbrace{H_S\otimes\mathbb{I}_E+\mathbb{I}_S\otimes H_E}^{H_0}+H_I\tag{\ref{eq:Htotal}}\\
    U(t,t_0) &= U_0(t,t_0)U_I(t,t_0)\tag{\ref{eq:Utotal}}
\end{align}
\textit{Recall:} $\ket{\psi(t)}_S=U(t)\ket{\psi}_S$
\begin{equation}
    \rho_S(t) = U(t)\rho_{S}U^\dagger(t)\label{eq:evolrhoU}
\end{equation}
If we now try to calculate the average (in $\mathcal{H}_T$) of an operator in the \textit{Schrödinger picture}, we can find a new interesting relation:
\begin{align*}
    \expval{O_S(t)} &= \Tr{\rho_{S}(t)O_S(t)}\\
    &= \Tr{U(t)\rho_{S}U^\dagger(t)O_S(t)}\\
    &= \Tr{U_0(t)U_I(t)\rho_{S}U_I^\dagger(t)U_0^\dagger(t)O_S(t)}\\
    &= \mathrm{Tr}\,[\underbrace{U_I(t)\rho_{S}U_I^\dagger(t)}_{\tilde{\rho}(t)}\,\underbrace{U_0^\dagger(t)O_S(t)U_0(t)}_{\tilde{O}(t)}]
\end{align*}
We can thus define in the interaction picture:
\begin{align}
    \tilde{\rho}(t)&=U_I(t)\rho_{S}U_I^\dagger(t)\tag{\ref{eq:rhoIP}}\\
    \tilde{O}(t) &= U_0^\dagger(t)O_S(t)U_0(t)\tag{\ref{eq:OIP}}
\end{align}
We can now reverse Eq.~(\ref{eq:evolrhoU}) and then use Eq.~(\ref{eq:Utotal}):
\begin{equation*}
    \rho_S = U^\dagger(t)\rho_S(t)U(t) = U_I^\dagger(t)U^\dagger_0(t)\rho_S(t)U_0(t)U_I(t)
\end{equation*}
And, by injecting this into Eq.~(\ref{eq:rhoIP}), we get
\begin{align}
    \tilde{\rho}(t) &= U_I(t)U_I^\dagger(t)U^\dagger_0(t)\rho_S(t)U_0(t)U_I(t)U^\dagger_I(t)\nonumber\\
    &= U^\dagger_0(t)\rho_S(t)U_0(t)\tag{\ref{eq:rhoIP}}
\end{align}
So that we find the results used in the main part of this work
\[\tilde{\rho}(t)=U_I(t)\rho_{S}U_I^\dagger(t)= U^\dagger_0(t)\rho_S(t)U_0(t)\]
% \section{Quantum Measurement Theory}\label{Measurement_Theory}

% The introduction about measurement theory in~\ref{intro} and this appendix 
% have been constructed from the work of \textit{K.Jacobs}~\cite{Jacobs} 
% and \textit{H.M.Wiseman}~\cite{WisemanThesis}.